\PassOptionsToPackage{table}{xcolor}
\documentclass[10pt,aspectratio=169]{beamer}
\newcommand{\LectureNumber}{08}
\input{course-theme.tex}
\title{Selection with if and else}
\subtitle{CSC103 Programming Fundamentals / Lecture 08}
\author{Dr. Muhammad Farhan}
\institute{Associate Professor of Computer Science\\Department of Computer Science\\COMSATS University Islamabad\\Sahiwal Campus}
\date{Fall 2026}
\begin{document}
\begin{frame}\titlepage\end{frame}

\begin{frame}[fragile,t]{The problem for this lecture}
\begin{columns}[T,onlytextwidth]\begin{column}{0.60\textwidth}
Classify an integer as negative, zero or positive. Draw a decision outline, then write an if-else chain.
\end{column}\begin{column}{0.35\textwidth}\small
Check the negative case first.
\end{column}\end{columns}
\note{Introduce the target problem before explaining the tools. Revisit it in the guided solution later.\par Syllabus reading: Deitel Ch 7 (as listed). CLO-2.}
\end{frame}

\begin{frame}[fragile,t]{Learning objectives}
\begin{columns}[T,onlytextwidth]\begin{column}{0.60\textwidth}
\begin{itemize}
\item Implement one-way and two-way selection
\item Build an ordered multi-branch decision
\item Test boundaries and nested conditions
\end{itemize}
\end{column}\begin{column}{0.35\textwidth}\small
CLO-2

The if statement\par Two-way selection\par Multiple selections\par Nesting and compound statements
\end{column}\end{columns}
\note{Explain how each objective supports the opening problem. The final questions check these objectives.\par Syllabus reading: Deitel Ch 7 (as listed). CLO-2.}
\end{frame}

\begin{frame}[fragile,t]{The if statement}
\begin{itemize}
\item An if condition must be Boolean.
\item The body executes only when the condition is true.
\item Braces clearly group all statements in the branch.
\end{itemize}
\vspace{1mm}\begin{block}{Example}
\begin{lstlisting}
if (balance < 0) {
  System.out.println("Overdrawn");
}
\end{lstlisting}
\end{block}
\note{Ask whether the program stops when the condition is false. It continues after the if block.\par Syllabus reading: Deitel Ch 7 (as listed). CLO-2.}
\end{frame}

\begin{frame}[fragile,t]{Two-way selection}
\begin{itemize}
\item An if-else chooses exactly one branch.
\item Statements after the structure execute after either branch completes normally.
\end{itemize}
\vspace{1mm}\begin{block}{Example}
\begin{lstlisting}
if (marks >= 50) {
  System.out.println("Pass");
} else {
  System.out.println("Fail");
}
\end{lstlisting}
\end{block}
\note{Trace 49, 50 and 51. Avoid implying these local examples redefine the syllabus assessment policy.\par Syllabus reading: Deitel Ch 7 (as listed). CLO-2.}
\end{frame}

\begin{frame}[fragile,t]{Multiple selections}
\begin{itemize}
\item An else-if chain selects the first matching branch.
\item Order overlapping thresholds from most restrictive to least restrictive.
\item Independent if statements can execute several bodies.
\end{itemize}
\vspace{1mm}\begin{block}{Example}
\begin{lstlisting}
if (score >= 80) {
  band = "High";
} else if (score >= 50) {
  band = "Middle";
} else { band = "Low"; }
\end{lstlisting}
\end{block}
\note{These are illustrative score bands, not official grades. Show how putting \textgreater{}=50 first makes \textgreater{}=80 unreachable in the chain.\par Syllabus reading: Deitel Ch 7 (as listed). CLO-2.}
\end{frame}

\begin{frame}[fragile,t]{Nesting and compound statements}
\begin{itemize}
\item A nested if expresses a decision inside another decision.
\item An else attaches to the nearest unmatched if.
\item Braces prevent misleading indentation.
\end{itemize}
\vspace{1mm}\begin{block}{Example}
\begin{lstlisting}
if (registered) {
  if (paid) {
    System.out.println("Admit");
  }
}
\end{lstlisting}
\end{block}
\note{Compare this nesting to registered \&\& paid. Both can express the same admission condition when no separate branch behaviour is needed.\par Syllabus reading: Deitel Ch 7 (as listed). CLO-2.}
\end{frame}

\begin{frame}[fragile,t]{Branch order determines the selected result}
\begin{itemize}
\item An else-if chain stops testing after its first matching condition.
\item Overlapping ranges need an order that gives the most specific case priority.
\item A final else covers every case not selected earlier.
\end{itemize}
\vspace{1mm}\begin{block}{Example}
\begin{lstlisting}
If >=50 comes before >=80, a score of 90
selects the >=50 branch immediately.
\end{lstlisting}
\end{block}
\note{Explain the mechanism using the displayed example. An else-if chain stops testing after its first matching condition. Overlapping ranges need an order that gives the most specific case priority. A final else covers every case not selected earlier.\par Syllabus reading: Deitel Ch 7 (as listed). CLO-2.}
\end{frame}

\begin{frame}[fragile,t]{Independent tests and exclusive alternatives}
\begin{itemize}
\item Separate if statements can execute more than one body.
\item An if-else chain expresses mutually exclusive outcomes.
\item For sign classification, exactly one of negative, zero and positive is appropriate.
\end{itemize}
\vspace{1mm}\begin{block}{Example}
\begin{lstlisting}
if (n < 0) ...
else if (n == 0) ...
else ...
\end{lstlisting}
\end{block}
\note{Explain the mechanism using the displayed example. Separate if statements can execute more than one body. An if-else chain expresses mutually exclusive outcomes. For sign classification, exactly one of negative, zero and positive is appropriate.\par Syllabus reading: Deitel Ch 7 (as listed). CLO-2.}
\end{frame}


\begin{frame}[fragile,t]{A decision selects one path}
\begin{center}\begin{tikzpicture}
\node[decision] (a) at (0,0) {mark \textgreater{}= 50?};
\node[alt] (b) at (-3.6,-1.65) {Print Pass};
\node[hot] (c) at (3.6,-1.65) {Print Fail};
\node[box] (d) at (0,-3.25) {Continue};
\draw[arr] (a.west) -| node[lab,pos=.25]{true} (b.north);\draw[arr] (a.east) -| node[lab,pos=.25]{false} (c.north);\draw[arr] (b.south) |- (d.west);\draw[arr] (c.south) |- (d.east);
\end{tikzpicture}\end{center}
\begin{block}{What to notice}Exactly one branch executes; both normal paths rejoin after the selection.\end{block}
\note{Trace each labelled connection. Exactly one branch executes; both normal paths rejoin after the selection.}
\end{frame}
\begin{frame}[fragile,t]{Key distinctions}
\small\renewcommand{\arraystretch}{1.5}
\rowcolors{2}{pale}{white}
\begin{tabularx}{\textwidth}{>{\raggedright\arraybackslash}X>{\raggedright\arraybackslash}X>{\raggedright\arraybackslash}X>{\raggedright\arraybackslash}X}
\rowcolor{navy}
\color{white}\bfseries Input n & \color{white}\bfseries n \textless{} 0 & \color{white}\bfseries n == 0 if reached & \color{white}\bfseries Chosen result\\
-1 & true & Not evaluated & Negative\\
0 & false & true & Zero\\
1 & false & false & Positive\\
\end{tabularx}
\vspace{3mm}\footnotesize These distinctions determine which operation or control structure fits the problem.
\note{\par Syllabus reading: Deitel Ch 7 (as listed). CLO-2.}
\end{frame}

\begin{frame}[fragile,t]{First example: execution}
\begin{lstlisting}
int score = 80;
String band;
if (score >= 80) { band = "High"; }
else if (score >= 50) { band = "Middle"; }
else { band = "Low"; }
System.out.println(band);
\end{lstlisting}
\note{This is the main body from Java\_Examples/Lecture08.java. The source file includes the complete class. Predict the result before the next slide.\par Syllabus reading: Deitel Ch 7 (as listed). CLO-2.}
\end{frame}

\begin{frame}[fragile,t]{First example: result}
\begin{columns}[T,onlytextwidth]\begin{column}{0.60\textwidth}
High\par The first condition is true.\par The remaining branches are skipped.
\end{column}\begin{column}{0.35\textwidth}\small
Multiple selections

Nesting and compound statements
\end{column}\end{columns}
\note{Expected output and interpretation: High
The first condition is true.
The remaining branches are skipped.\par Syllabus reading: Deitel Ch 7 (as listed). CLO-2.}
\end{frame}

\begin{frame}[fragile,t]{Guided practice}
\begin{columns}[T,onlytextwidth]\begin{column}{0.60\textwidth}
Classify an integer as negative, zero or positive. Draw a decision outline, then write an if-else chain.
\end{column}\begin{column}{0.35\textwidth}\small
Attempt the solution before continuing.

Use the definitions and examples from this lecture.
\end{column}\end{columns}
\note{Give students 8 minutes to attempt the problem. The following slides provide a complete solution. Original solution guidance: if (n \textless{} 0) ... else if (n == 0) ... else ...
Test -1, 0 and 1. Exactly one label must appear.\par Syllabus reading: Deitel Ch 7 (as listed). CLO-2.}
\end{frame}

\begin{frame}[fragile,t]{Solution strategy}
\begin{columns}[T,onlytextwidth]\begin{column}{0.60\textwidth}
\begin{itemize}
\item Check the negative case first.
\item If the number is not negative, test equality to zero.
\item Every remaining integer is positive.
\end{itemize}
\end{column}\begin{column}{0.35\textwidth}\small
The next program uses fixed inputs so its result can be reproduced.
\end{column}\end{columns}
\note{Discuss why each step is needed. State the input assumptions before executing the program.\par Syllabus reading: Deitel Ch 7 (as listed). CLO-2.}
\end{frame}

\begin{frame}[fragile,t]{Complete solution}
\begin{lstlisting}
public class Solution08 {
  public static void main(String[] args) throws Exception {
    int n = 0;
    if (n < 0) {
      System.out.println("Negative");
    } else if (n == 0) {
      System.out.println("Zero");
    } else {
      System.out.println("Positive");
    }
  }
}
\end{lstlisting}
\note{Complete runnable file: Java\_Examples/Solution08.java. Inputs are fixed in the program.  \par Syllabus reading: Deitel Ch 7 (as listed). CLO-2.}
\end{frame}

\begin{frame}[fragile,t]{Execution trace}
\small\renewcommand{\arraystretch}{1.5}
\rowcolors{2}{pale}{white}
\begin{tabularx}{\textwidth}{>{\raggedright\arraybackslash}X>{\raggedright\arraybackslash}X>{\raggedright\arraybackslash}X}
\rowcolor{navy}
\color{white}\bfseries Execution point & \color{white}\bfseries Condition or action & \color{white}\bfseries Result\\
First test & 0 \textless{} 0 & false\\
Second test & 0 == 0 & true\\
Selected body & Print Zero & Zero\\
Remaining branch & Skip else & No extra output\\
\end{tabularx}
\vspace{3mm}\footnotesize Follow the changing state in execution order.
\note{Manually derived trace for the complete solution. Check the negative case first. If the number is not negative, test equality to zero. Every remaining integer is positive.\par Syllabus reading: Deitel Ch 7 (as listed). CLO-2.}
\end{frame}

\begin{frame}[fragile,t]{Solution result}
\begin{columns}[T,onlytextwidth]\begin{column}{0.60\textwidth}
Zero
\end{column}\begin{column}{0.35\textwidth}\small
Why this result follows

Every remaining integer is positive.
\end{column}\end{columns}
\note{Expected result: Zero\par Syllabus reading: Deitel Ch 7 (as listed). CLO-2.}
\end{frame}

\begin{frame}[fragile,t]{A common incorrect approach}
if (marks \textgreater{}= 50);\par \{ System.out.println("Pass"); \}
\vspace{1mm}\begin{block}{Example}
\begin{lstlisting}
The semicolon creates an empty if body. Remove it so that the block belongs to the if.
\end{lstlisting}
\end{block}
\note{Ask students to predict the failure before discussing the explanation. The right side gives the correction.\par Syllabus reading: Deitel Ch 7 (as listed). CLO-2.}
\end{frame}

\begin{frame}[fragile,t]{Boundary and failure checks}
\begin{columns}[T,onlytextwidth]\begin{column}{0.60\textwidth}
Check the stated input assumptions.

Write a shipping-charge decision for three weight ranges. State your rules and test each boundary.
\end{column}\begin{column}{0.35\textwidth}\small
if (n \textless{} 0) ... else if (n == 0) ... else ...\par Test -1, 0 and 1. Exactly one label must appear.
\end{column}\end{columns}
\note{Use the specification to derive each expected result. Exercise solution: if (n \textless{} 0) ... else if (n == 0) ... else ...
Test -1, 0 and 1. Exactly one label must appear.\par Syllabus reading: Deitel Ch 7 (as listed). CLO-2.}
\end{frame}

\begin{frame}[fragile,t]{Check your understanding}
\begin{columns}[T,onlytextwidth]\begin{column}{0.60\textwidth}
How many branches can an else-if chain select? Why test 49, 50 and 51?
\end{column}\begin{column}{0.35\textwidth}\small
Write a prediction and explain the rule that supports it.
\end{column}\end{columns}
\note{Answer: At most one. Those cases exercise both sides of the 50 boundary and equality itself.\par Syllabus reading: Deitel Ch 7 (as listed). CLO-2.}
\end{frame}

\begin{frame}[fragile,t]{Answer and explanation}
\begin{columns}[T,onlytextwidth]\begin{column}{0.60\textwidth}
At most one. Those cases exercise both sides of the 50 boundary and equality itself.
\end{column}\begin{column}{0.35\textwidth}\small
The semicolon creates an empty if body. Remove it so that the block belongs to the if.
\end{column}\end{columns}
\note{Reconnect the answer to the relevant definition rather than treating it as a fact to memorise.\par Syllabus reading: Deitel Ch 7 (as listed). CLO-2.}
\end{frame}


\begin{frame}[fragile,t]{Compare cases and predict outcomes}
\small\renewcommand{\arraystretch}{1.5}\rowcolors{2}{pale}{white}
\begin{tabularx}{\textwidth}{>{\raggedright\arraybackslash}X>{\raggedright\arraybackslash}X>{\raggedright\arraybackslash}X}\rowcolor{navy}
\color{white}\bfseries Mark & \color{white}\bfseries Condition & \color{white}\bfseries Printed result\\
49 & false & Fail\\
50 & true & Pass\\
51 & true & Pass\\
\end{tabularx}\par\vspace{3mm}\begin{exampleblock}{Discuss}Explain each expected result before running the example. Which case would reveal a wrong assumption?\end{exampleblock}
\note{Read the first two columns and invite predictions before discussing the outcome.}
\end{frame}

\begin{frame}[fragile,t]{Apply the idea}
\begin{alertblock}{Independent practice}Write a grade band selection: 80 or above is High, 50 through 79 is Pass, and below 50 is Fail. Explain the order of tests.\end{alertblock}\vspace{3mm}\begin{enumerate}\item Work through the input by hand.\item Write or correct the program.\item Justify the expected result.\end{enumerate}\note{Allow 3–5 minutes, then compare answers in pairs. Write a grade band selection: 80 or above is High, 50 through 79 is Pass, and below 50 is Fail. Explain the order of tests.}
\end{frame}

\begin{frame}[fragile,t]{Solution and reasoning}
\begin{lstlisting}
if (mark >= 80) {
  System.out.println("High");
} else if (mark >= 50) {
  System.out.println("Pass");
} else {
  System.out.println("Fail");
}
\end{lstlisting}\vspace{1mm}\small\begin{itemize}
\item The narrower, higher threshold must be tested first.
\item After mark \textgreater{}= 80 is false, mark \textgreater{}= 50 means 50 through 79.
\item The example assumes a valid mark; range validation is a separate step.
\end{itemize}\note{Answer: The narrower, higher threshold must be tested first. After mark \textgreater{}= 80 is false, mark \textgreater{}= 50 means 50 through 79. The example assumes a valid mark; range validation is a separate step.}
\end{frame}

\begin{frame}[fragile,t]{Overtime payroll}
\begin{itemize}
\item The first 40 hours use the regular rate. Only excess hours use the 1.5 multiplier.
\item At exactly 40 hours, overtime hours are zero.
\item This is the fictional business rule supplied by the exercise, not an employment{-}law rule.
\end{itemize}
\vspace{3mm}\footnotesize\textcolor{accent}{Source: supplied ISB campus slides. See speaker notes and ISB\_Source\_Map.md.}
\note{Adapted from the supplied ISB campus folder: Lecture{-}7.pptx, slides 24, 25. Adapted explanation and runnable implementation; sample inputs and traces added for this course.}
\end{frame}

\begin{frame}[fragile,t]{Worked problem: Overtime payroll}
\begin{block}{Task}Compute pay for 45 hours at 20 units per hour using the source overtime rule.\end{block}
\small Input: Values are fixed in the program for a reproducible demonstration.\par\vspace{3mm}\begin{exampleblock}{Before running}Predict the output and identify the variables that change.\end{exampleblock}
\note{Adapted from the supplied ISB campus folder: Lecture{-}7.pptx, slides 24, 25. Adapted explanation and runnable implementation; sample inputs and traces added for this course.}
\end{frame}

\begin{frame}[fragile,t]{Complete ISB example (1/1)}
\begin{lstlisting}
public class IsbLecture08 {
  public static void main(String[] args) throws Exception {
    double hours = 45, rate = 20;
    double pay;
    if (hours > 40) {
      pay = 40 * rate + (hours - 40) * rate * 1.5;
    } else {
      pay = hours * rate;
    }
    System.out.printf(java.util.Locale.US, "Pay: %.2f%n", pay);
  }

}
\end{lstlisting}
\note{Adapted from the supplied ISB campus folder: Lecture{-}7.pptx, slides 24, 25. Adapted explanation and runnable implementation; sample inputs and traces added for this course. Complete file: ISB\_Java\_Examples/IsbLecture08.java.}
\end{frame}

\begin{frame}[fragile,t]{Worked trace and expected result}
\small\renewcommand{\arraystretch}{1.4}\rowcolors{2}{pale}{white}
\begin{tabularx}{\textwidth}{>{\raggedright\arraybackslash}X>{\raggedright\arraybackslash}X>{\raggedright\arraybackslash}X}\rowcolor{navy}
\color{white}\bfseries Hours & \color{white}\bfseries Regular + overtime pay & \color{white}\bfseries Total\\
39 & 780 + 0 & 780\\
40 & 800 + 0 & 800\\
45 & 800 + 150 & 950\\
\end{tabularx}\par\vspace{2mm}
\begin{block}{Expected console output}\ttfamily\footnotesize Pay: 950.00\end{block}
\note{Adapted from the supplied ISB campus folder: Lecture{-}7.pptx, slides 24, 25. Adapted explanation and runnable implementation; sample inputs and traces added for this course. Trace and expected values independently worked for this edition.}
\end{frame}

\begin{frame}[fragile,t]{Extend the ISB example}
\begin{alertblock}{Practice}What is wrong with multiplying all 45 hours by 1.5 times the rate?\end{alertblock}\vspace{3mm}\begin{exampleblock}{Explain your reasoning}It applies the overtime premium to regular hours. That incorrect formula gives 1350 instead of 950.\end{exampleblock}
\note{Adapted from the supplied ISB campus folder: Lecture{-}7.pptx, slides 24, 25. Adapted explanation and runnable implementation; sample inputs and traces added for this course. Answer: It applies the overtime premium to regular hours. That incorrect formula gives 1350 instead of 950.}
\end{frame}
\begin{frame}[fragile,t]{Lecture summary}
\begin{columns}[T,onlytextwidth]\begin{column}{0.60\textwidth}
\begin{itemize}
\item one-way and two-way selection
\item an ordered multi-branch decision
\item boundaries and nested conditions
\end{itemize}
\end{column}\begin{column}{0.35\textwidth}\small
\begin{itemize}
\item Check the negative case first.
\item If the number is not negative, test equality to zero.
\item Every remaining integer is positive.
\end{itemize}
\end{column}\end{columns}
\note{Ask students to give one concrete example for the concepts listed. Use the worked solution to connect the topics.\par Syllabus reading: Deitel Ch 7 (as listed). CLO-2.}
\end{frame}

\begin{frame}[fragile,t]{After-class practice}
\begin{columns}[T,onlytextwidth]\begin{column}{0.60\textwidth}
Write a shipping-charge decision for three weight ranges. State your rules and test each boundary.
\end{column}\begin{column}{0.35\textwidth}\small
Syllabus reading\par Deitel Ch 7 (as listed)

Next lecture: Short-circuiting and switch
\end{column}\end{columns}
\note{Reading labels follow the supplied outline. Check topic headings against the available textbook edition. Require a program and expected results for the practice task.\par Syllabus reading: Deitel Ch 7 (as listed). CLO-2.}
\end{frame}
\end{document}
